\section{Limitations and Threats to Validity}
\label{sec:limitations}

We list the limitations explicitly because they bound the external validity of every headline number. None is fatal; each carries an associated mitigation or is being addressed in the cross-application work of Section~\ref{sec:cross-app}.

\subsection{Dataset realism}

\begin{enumerate}
\item \textbf{Trace sampling is 100\%.} Production typically samples traces at 1--10\% head-sampling or tail-on-error. Models trained on densely-sampled traces may over-rely on span-fan-out features that are partly missing under production sampling. We chose dataset density over sampling realism and disclose the gap.

\item \textbf{Single-cluster, single-region deployment.} No cross-availability-zone failures, no inter-region partitions, no multi-cluster failover behavior. Cross-region scenarios deferred.

\item \textbf{Synthetic traffic.} The load-generator emits a deterministic basket of routes with no daily, weekly, or business-hour pattern. Real production traffic carries seasonality that affects baseline metrics; the absence of seasonality narrows the realistic noise distribution.

\item \textbf{Single fault per run.} Real production incidents often have correlated multi-fault cascades (a deploy-induced latency plus a cache eviction plus a downstream slow query, all firing within minutes). Our multi-fault scenarios are limited; the OTel Demo extension is the place we are turning this from a limitation into a contribution.

\item \textbf{Shadow Jira, not engineer-authored.} Every ticket in the memory corpus was generated from a scenario YAML template and humanized by an LLM. Real engineer-written tickets carry inconsistencies, typos, out-of-band context (Slack threads, postmortem links, manager pressure) that synthetic tickets do not.

\item \textbf{No alert-fatigue baseline.} Real on-call queues have background false-positive noise from alert rules that pre-date the period under study. Our baseline-normal family emits zero alerts.

\item \textbf{No GPU-served services.} Real recommendation and ad services often use GPU-backed inference; their failures (CUDA OOM, model cold-start, quota throttling) are absent from our dataset.
\end{enumerate}

\subsection{Methodological constraints}

\begin{enumerate}
\setcounter{enumi}{7}
\item \textbf{In-distribution evaluation dominates.} The 1{,}008-window test split has same-family windows in train. The G8 LOFO analysis is the only out-of-distribution evaluation; it shows reasonable generalization but only at the family-shift level. Application-shift generalization is what Section~\ref{sec:cross-app} addresses; until that collection finishes, the application-shift claim is unproven.

\item \textbf{Single dataset.} Every retrieval claim is on one V2 humanized corpus of 347 tickets. A real Jira corpus could differ substantially in noise, length, vocabulary, and structure --- which is why we shaped the V2 humanizer against TAWOS and report cross-corpus retrieval as future work.

\item \textbf{Distractor analysis is a simulation.} The G6 robustness sweep simulates distractor injection at the prediction level rather than re-fitting the bi-encoder + SPLADE + KG per distractor ratio. The numbers are a lower bound on degradation; real re-fits could be modestly worse.

\item \textbf{The $+1.7$-point Hit@1 lift over BiEncoder is not statistically significant.} The cascade ties BiEncoder on top-1 within the 95\% CI. The big wins are on Hit@5 (broader coverage), MRR, and novelty --- not on top-1 precision.

\item \textbf{Agent novelty recall is bounded.} Even after running the agent on all 1{,}008 windows, agent-attributable novelty recall plateaus around $35\%$. The G7 jump to $79\%$ comes from the learned signal, not the agent. If the learned signal degrades under genuine distribution shift, novel recall would fall back toward $35\%$.

\item \textbf{Triage performance attribution.} HGB alone achieves strict PR-AUC $= 0.9998$ on this dataset. The cascade's triage contribution is on borderline windows ($+3.2$ points inclusive PR-AUC). If HGB's numeric-feature signal degrades on new data (different services, different metric distributions), the cascade's triage degrades with it.

\item \textbf{Cascade is composition-fragile.} As G3 showed, a large standalone gain on one component can break cascade integration through implicit score-scale assumptions. Future cascade extensions must address score normalization explicitly.
\end{enumerate}

\subsection{Non-claims --- what we deliberately do not claim}

\begin{itemize}
\item Memory augmentation improves anomaly detection. HGB alone already saturates this dataset.
\item Our system outperforms commercial observability tools. We have no commercial-tool comparison; out of scope.
\item The system is production-ready. It is a research prototype on a synthetic-but-realistic corpus.
\item We solve cold-start. With zero past tickets, retrieval is $0\%$ by definition.
\item The humanized Jira corpus is indistinguishable from real Jira. We claim engineer-voice realism on the dimensions that matter for retrieval (length, multi-channel evidence, log signatures); we do not claim parity.
\end{itemize}
